\documentclass[12pt,a4paper]{article}

% ── Fonts and encoding (LuaLaTeX) ─────────────────────────────────────────────
\usepackage{fontspec}
\usepackage{unicode-math}
\setmainfont{Linux Libertine O}
\setmathfont{TeX Gyre Pagella Math}
\newfontfamily\shavfont{Linux Libertine O}
\newfontfamily\igprimfont{Everson Mono}
\setmonofont[Scale=0.85]{Everson Mono}

% ── Layout ────────────────────────────────────────────────────────────────────
\usepackage[top=1in, bottom=1in, left=1.1in, right=1.1in, headheight=15pt]{geometry}
\usepackage{microtype}
\usepackage{parskip}
\usepackage[english]{babel}

% ── Headers ───────────────────────────────────────────────────────────────────
\usepackage{fancyhdr}
\pagestyle{fancy}
\fancyhf{}
\fancyhead[L]{\leftmark}
\fancyhead[R]{\thepage}
\renewcommand{\headrulewidth}{0.4pt}

% ── Section styling ───────────────────────────────────────────────────────────
\usepackage{titlesec}
\titleformat{\section}{\Large\bfseries}{\thesection}{1em}{}
\titleformat{\subsection}{\large\bfseries}{\thesubsection}{1em}{}
\titleformat{\subsubsection}{\normalsize\bfseries}{\thesubsubsection}{1em}{}

% ── Lists ─────────────────────────────────────────────────────────────────────
\usepackage{enumitem}
\setlist[itemize]{topsep=0.5ex, itemsep=0.25ex, parsep=0pt}
\setlist[enumerate]{topsep=0.5ex, itemsep=0.25ex, parsep=0pt}

% ── Math and tables ───────────────────────────────────────────────────────────
\usepackage{amsmath,amsthm}
\usepackage{mathtools}
\usepackage{booktabs}
\usepackage{array}
\usepackage{tabularx}
\usepackage[table]{xcolor}
\renewcommand{\arraystretch}{1.3}

% ── Code listings ─────────────────────────────────────────────────────────────
\usepackage{listings}
\definecolor{codebg}{RGB}{248,248,248}
\definecolor{codeframe}{RGB}{200,200,200}
\definecolor{codekeyword}{RGB}{0,0,180}
\definecolor{codecomment}{RGB}{0,128,0}
\lstset{
  basicstyle=\ttfamily\footnotesize,
  backgroundcolor=\color{codebg},
  frame=single,
  rulecolor=\color{codeframe},
  keywordstyle=\color{codekeyword}\bfseries,
  commentstyle=\color{codecomment}\itshape,
  breaklines=true,
  captionpos=b,
  morekeywords={axiom,def,theorem,noncomputable,structure,where,namespace,
                open,import,section,end,fun,exact,intro,simp,rfl,decide,
                constructor,obtain,refine,cases,by,have,exact,use,rw},
}

% ── Cross-references ──────────────────────────────────────────────────────────
\usepackage{hyperref}
\usepackage{cleveref}
\hypersetup{colorlinks=true, linkcolor=blue, urlcolor=cyan, citecolor=blue}

% ── Theorem environments ──────────────────────────────────────────────────────
\newtheorem{theorem}{Theorem}[section]
\newtheorem{lemma}[theorem]{Lemma}
\newtheorem{proposition}[theorem]{Proposition}
\newtheorem{corollary}[theorem]{Corollary}
\newtheorem{definition}[theorem]{Definition}
\newtheorem{axiom_env}[theorem]{Axiom}
\newtheorem{remark}[theorem]{Remark}

% ── Custom notation ───────────────────────────────────────────────────────────
\newcommand{\BFour}{\mathbf{B}_4}
\newcommand{\WH}{\mathrm{WH}}
\newcommand{\SIC}{\mathrm{SIC}}
\newcommand{\Kd}{K_d}
\newcommand{\Fd}{F_d}
\newcommand{\md}{m_d}
\newcommand{\frob}{\mu \circ \delta = \mathrm{id}}
\newcommand{\Lean}{\textsc{Lean\,4}}

% ─────────────────────────────────────────────────────────────────────────────

\begin{document}

\title{\textbf{SIC-POVMs, the Stark Conjecture, and Hilbert's Twelfth Problem:\
A Lean\,4 Formalization via Paraconsistent Belnap Multilattices}}

\author{C.\ L.\ Mills (Lando\,$\otimes$\,$\odot$-operator)%
\thanks{Thanks are owed to Harry T.\ Larson; cf.\ \cite{Larson1961}.}}

\date{June 2026}

\maketitle

\begin{abstract}
We present a \Lean{} formalization establishing a three-level structural
identification: \emph{(i)}~the Belnap multilattice axioms for
Weyl--Heisenberg (WH) covariant symmetric informationally complete
positive operator-valued measures (SIC-POVMs) in dimension $d = 2^n$;
\emph{(ii)}~the Zauner conjecture for those dimensions; and
\emph{(iii)}~the mixed-signature Stark conjecture for the ray class fields
$\Kd = \mathbb{Q}(\sqrt{d(d-2)})$ --- a special case of Hilbert's Twelfth
Problem for real quadratic fields. The core equivalence
\texttt{hilbert\_embedding\_equiv\_zauner} is proved by \texttt{rfl}: the
Hilbert-space embedding of the Belnap multilattice into $\mathbb{C}^{2^n}$
and the Zauner conjecture at $d = 2^n$ are definitionally identical
propositions. The Belnap structural skeleton---orbit cardinality $4^n$,
Frobenius closure $\frob$, join-equiangularity, and a structural Born
rule---is proved unconditionally with zero \texttt{sorry}s.
The remaining open content is precisely axiomatized: three named axioms
carry the exact arithmetic geometry of Stark units acting on WH frames.
A proof of the mixed-signature Stark conjecture closes all three levels
simultaneously.
\end{abstract}

\tableofcontents
\newpage

%──────────────────────────────────────────────────────────────────────────────
\section{Introduction}
\label{sec:intro}
%──────────────────────────────────────────────────────────────────────────────

Symmetric informationally complete positive operator-valued measures
(SIC-POVMs) are minimal tomographically complete quantum measurements: $d^2$
rank-1 projectors in $\mathbb{C}^d$ with constant pairwise Hilbert--Schmidt
inner products.  Their existence in every finite dimension is widely
conjectured (Zauner's conjecture, 1999 \cite{Zauner1999}) and verified
numerically through $d = 2406$, yet no general proof exists.

Appleby, Bengtsson, Flammia, and collaborators \cite{Appleby2013,Appleby2017}
discovered that SIC-POVMs have unexpected arithmetic structure: the fiducial
vector coordinates lie in specific ray class fields of real quadratic fields.
The base field for dimension $d$ is $F_d = \mathbb{Q}(\sqrt{d(d-2)})$, and
the coordinates of the fiducial generate the ray class field $\Kd$ of $F_d$.
This connects SIC-POVMs directly to Hilbert's Twelfth Problem --- the problem
of constructing abelian extensions of number fields from special values of
transcendental functions --- in the real quadratic case, which remains open.

The present work formalizes these connections in \Lean{} \cite{Lean4}
(v4.28.0, paraconsistent kernel fork \cite{p4rakernel}) using:
\begin{itemize}
  \item \texttt{Imscribing.Millennium.SIC\_POVM\_Stark}: Weyl--Heisenberg
    operators in $\mathbb{C}^d$, the SIC-POVM definition, Stark arithmetic,
    and the conditional existence theorem (\S\ref{sec:stark}).
  \item \texttt{Imscribing.Paraconsistent.Shor.BelnapWHMultilattice}: the
    Belnap $\BFour$ multilattice, the product-lattice orbit theorem
    ($2^n$ proved), and the multilattice axioms for $4^n$ (\S\ref{sec:belnap}).
  \item \texttt{Imscribing.Millennium.ZaunerEmbeddingEquivalence}: the main
    equivalence and its corollaries (\S\ref{sec:equiv}).
\end{itemize}

The structural picture is:

\[
\underbrace{\text{Belnap multilattice}}_{\text{sorry-free}}
\;\xrightarrow{\;\texttt{rfl}\;}
\;[\,\texttt{HilbertEmbeddingExists} \;\leftrightarrow\; \texttt{ZaunerConjectureAtPow2}\,]
\;\xrightarrow{\;\text{axiom: \texttt{stark\_equivalence}}\;}
\;\texttt{SICPOVM\_Exists}\,d
\;\xrightarrow{\;\text{axioms: \texttt{equiangular} + \texttt{norm}}\;}
\;\texttt{MixedSignatureStarkConjecture}
\;\xrightarrow{\;\text{axiom: \texttt{zauner\_correspondence}}\;}
\;\text{Stark units in }\Kd \;=\; \text{Hilbert\,12 for }\mathbb{Q}(\sqrt{d(d-2)})
\]

Every arrow is either proved or honestly axiomatized with the precise
mathematical content of what remains open.

%──────────────────────────────────────────────────────────────────────────────
\section{Weyl--Heisenberg Operators and the SIC-POVM Definition}
\label{sec:wh}
%──────────────────────────────────────────────────────────────────────────────

\subsection{The Weyl--Heisenberg group in dimension $d$}

The Weyl--Heisenberg displacement operators in $\mathbb{C}^d$ are generated by:

\begin{definition}[Shift and phase operators]
Let $\omega_d = e^{2\pi i/d}$.  Define:
\begin{align*}
  (X_d\,v)(k) &= v(k-1 \bmod d) & (\text{shift}) \
  (Z_d\,v)(k) &= \omega_d^k \cdot v(k) & (\text{phase})
\end{align*}
The displacement operator is $D_{a,b,t} = \omega_d^t X_d^a Z_d^b$.
\end{definition}

In \Lean{} (\texttt{SIC\_POVM\_Stark.lean}, lines 26--43):

\begin{lstlisting}
def omega_d (d : N) : C := exp (2 * Real.pi * Complex.I / d)
def X_d (d : N) (v : Fin d -> C) (k : Fin d) : C :=
  v <<(k.val - 1) % d, Nat.mod_lt _ k.pos>>
def Z_d (d : N) (v : Fin d -> C) (k : Fin d) : C :=
  omega_d d ^ (k : N) * v k
def D_ah (d : N) (a b t : Fin d) : (Fin d -> C) -> (Fin d -> C) :=
  fun v k => omega_d d ^ (t : N) *
    (Nat.iterate (X_d d) (a : N) (Nat.iterate (Z_d d) (b : N) v)) k
\end{lstlisting}

These are genuine displacement operators, not placeholder types.

\subsection{The SIC-POVM condition}

\begin{definition}[IsSICPOVM]
A vector $\phi \in \mathbb{C}^d$ is a \emph{SIC-POVM fiducial} if:
\begin{enumerate}[label=(\roman*)]
  \item $\langle \phi, \phi \rangle = d$ \quad (norm condition), and
  \item $|\langle \phi, D_{a,b,0}\,\phi \rangle| = 1$ for all $(a,b) \neq (0,0)$ \quad (equiangularity).
\end{enumerate}
\end{definition}

\begin{lstlisting}
structure IsSICPOVM (d : N) [NeZero d] (fiducial : Fin d -> C) : Prop where
  norm_eq    : wh_normSq d fiducial = (d : R)
  equiangular : forall (a b : Fin d), (a, b) != (0, 0) ->
    ||wh_inner d fiducial (D_ah d a b 0 fiducial)|| = 1
\end{lstlisting}

The Frobenius inner product is $\langle v, w \rangle = \sum_k v_k \overline{w_k}$,
implemented as \texttt{wh\_inner}.  The condition $|\langle\phi, D\phi\rangle|=1$
(not $1/(d+1)$) is the un-normalized formulation; the normalized fiducial
\texttt{normalize\_fiducial} satisfies the standard $1/(d+1)$ overlap.

%──────────────────────────────────────────────────────────────────────────────
\section{Arithmetic Geometry: the Stark Conjecture}
\label{sec:stark}
%──────────────────────────────────────────────────────────────────────────────

\subsection{The discriminant and base field}

\begin{definition}[Base field]
For $d \geq 2$, let $m_d = d(d-2)$ and $F_d = \mathbb{Q}(\sqrt{m_d})$.
\end{definition}

The definition \texttt{m\_d (d : N) : Z := (d : Z) * ((d : Z) - 2)} is
exactly the discriminant of $F_d$.  The dimension $d$ \emph{selects} the
arithmetic field: the SIC-POVM problem in dimension $d$ is the Stark problem
over $F_d$.  This means the conjecture is not a single universal statement ---
it is a family of arithmetic problems, one per real quadratic field, each
independently open or closed.

\begin{remark}
For $d=3$: $m_3 = 3$, $F_3 = \mathbb{Q}(\sqrt{3})$ --- the Hecke field of
the simplest equilateral SIC.  For $d=4$: $m_4 = 8$, $F_4 = \mathbb{Q}(\sqrt{2})$.
Both are known to support SIC-POVMs; the Stark units are explicitly known in
these cases.
\end{remark}

\subsection{The mixed-signature Stark conjecture}

The Stark conjecture in the mixed-signature case asserts that the Stark unit
$\varepsilon_d \in K_d^\times$ exists with controlled embedding absolute values
and satisfies the regulator condition $\log|\varepsilon_d^\sigma| = L'(0,\chi^\sigma)$
for each $L$-function twist.

\begin{axiom_env}[MixedSignatureStarkConjecture]
\label{ax:stark}
For $d \geq 2$ with $m_d$ non-square:
\[
\forall\, i,j \in \{0,\ldots,d-1\}:\quad
\|\sigma_i(\varepsilon_d)\| \leq \tfrac{1}{d} + 1
\quad\text{and}\quad
\forall\, \tau \in \mathrm{Gal}(K_d/F_d):\; \sigma_j(\tau\cdot\varepsilon_d) \neq 0
\]
\end{axiom_env}

This is an open problem in analytic number theory.  The \Lean{}
formalization marks it as \texttt{def MixedSignatureStarkConjecture}, not
\texttt{axiom} --- it is a defined proposition whose truth value is unknown.

\subsection{Constructing the fiducial from the Stark unit}

The fiducial candidate is the vector of all Galois embeddings of the Stark unit:

\[
v_d(k) = \sigma_k(\varepsilon_d),\quad k = 0,\ldots,d-1
\]

\begin{lstlisting}
def fiducial_from_stark (d : N) (hd : 2 <= d) (hns : not IsSquare (m_d d)) :
    Fin d -> C :=
  fun k => (Embeddings d hd hns k) (StarkUnit d hd hns)

def normalize_fiducial ... :=
  fun k => (Real.sqrt (d : R))^(-1) * fiducial_from_stark d hd hns k
\end{lstlisting}

\subsection{The Galois--Zauner correspondence}

\begin{axiom_env}[zauner\_correspondence]
\label{ax:zauner-corr}
The absolute value of the WH orbit inner product is controlled by the
order-3 Zauner automorphism $\zeta_d \in \mathrm{Gal}(K_d/F_d)$ acting
on the Stark unit:
\[
\big|\langle v_d,\, D_{a,b}\,v_d\rangle\big|
= \big|\langle v_d, v_d\rangle\big| \cdot
\big|\overline{\sigma_0(\zeta_d \cdot \varepsilon_d)}\big|
\]
\end{axiom_env}

This axiom captures the Appleby--Bengtsson result: the SIC overlap is an
algebraic function of the Stark unit under the Zauner automorphism.  It is
formalized as a named \texttt{axiom} with the exact mathematical content
stated in the type.

\subsection{Honest gap markers}

Two further axioms carry the open arithmetic geometry:

\begin{axiom_env}[equiangular\_from\_stark\_axiom, norm\_of\_normalized\_axiom]
\label{ax:gap}
Given the mixed-signature Stark conjecture~\ref{ax:stark}, the normalized
fiducial satisfies:
\begin{itemize}
  \item $|\langle \tilde{v}_d, D_{a,b}\,\tilde{v}_d\rangle| = 1$ for all $(a,b)\neq(0,0)$
  \item $\langle\tilde{v}_d, \tilde{v}_d\rangle_{\mathbb{R}} = d$
\end{itemize}
\emph{Gap:} the full arithmetic geometry of Stark units acting on WH frames --- the
translation from the Stark embedding bounds to the exact equiangularity and norm
conditions --- is the open mathematical content.
\end{axiom_env}

Both axioms carry the same comment in the source:
\begin{quotation}
\noindent\textit{``the gap between the Stark conjecture and these norm statements
is the open content --- it requires the full arithmetic geometry of Stark units
acting on WH frames.''}
\end{quotation}

This is not evasion.  It is the precise location of what remains to be done.

\subsection{The conditional existence theorem}

\begin{theorem}[sic\_povm\_exists\_via\_stark]
\label{thm:stark-sic}
Assume the mixed-signature Stark conjecture and the two gap axioms
(\ref{ax:gap}).  Then $\texttt{SICPOVM\_Exists}(d)$ holds for all $d \geq 2$
with $m_d$ non-square.
\end{theorem}

\begin{lstlisting}
theorem sic_povm_exists_via_stark
    (d : N) [NeZero d] (hd : 2 <= d) (hns : not IsSquare (m_d d))
    (sc : MixedSignatureStarkConjecture d hd hns) :
    SICPOVM_Exists d := by
  use normalize_fiducial d hd hns
  exact { norm_eq := norm_of_normalized d hd hns sc,
          equiangular := equiangular_from_stark d hd hns sc }
\end{lstlisting}

%──────────────────────────────────────────────────────────────────────────────
\section{The Belnap Multilattice}
\label{sec:belnap}
%──────────────────────────────────────────────────────────────────────────────

\subsection{The Belnap four-valued lattice $\BFour$}

The Belnap--Dunn four-valued logic \cite{Belnap1977} has truth values
$\{N, T, F, B\}$ (neither, true, false, both) ordered by the bilattice
structure.  The key property for our purposes:
\[
\neg B = B \qquad (\texttt{bnot B = B})
\]
This is the \emph{dialetheic fixed point}: $B$ absorbs negation.

\subsection{The $n$-qubit Belnap state and WH action}

An $n$-qubit Belnap state is a pair $(\text{val}, \text{phase}) \in \BFour \times \mathbb{Z}_2$.
The all-$B$ all-zero state $B^{\otimes n}$ is the fiducial.  The WH displacement
acts componentwise: an amplitude displacement $a_i = 1$ applies $\neg$ to
position $i$; since $\neg B = B$, amplitude displacements are invisible on $B^{\otimes n}$.

\begin{theorem}[whOrbit\_card\_eq\_pow2]
\label{thm:orbit-2n}
The product-lattice WH orbit of $B^{\otimes n}$ has cardinality $2^n$:
\[
|\mathcal{O}_{B^{\otimes n}}| = 2^n
\]
\end{theorem}

\begin{proof}
Amplitude displacements fix $B^{\otimes n}$ (since $\neg B = B$).  The displaced
states differ only in the phase word $p \in (\mathbb{Z}_2)^n$.  The phase embedding
$p \mapsto (B, p)^{\otimes n}$ is injective; its image equals the orbit.
Hence $|\mathcal{O}| = |(\mathbb{Z}_2)^n| = 2^n$.  \textbf{(Proved in \Lean{}, 0 sorries.)}
\end{proof}

\begin{corollary}[product\_lattice\_orbit\_is\_insufficient]
\label{cor:gap}
For $n \geq 1$: $|\mathcal{O}_{B^{\otimes n}}| = 2^n < 4^n = d^2$.
The product-lattice orbit is insufficient for a SIC-POVM.
\end{corollary}

This is the \emph{orbit gap} --- the structural gap that the multilattice
axioms close.

\subsection{The multilattice axioms}

An extended type $\texttt{BelnapML}(n)$ is postulated to support amplitude-distinguishable
displacements.  Four axioms specify its properties:

\begin{axiom_env}[Ax-PROJ]
The multilattice fiducial projects to $B^{\otimes n}$ in the product lattice.
\end{axiom_env}

\begin{axiom_env}[Ax-FREE]
The WH action on the multilattice fiducial produces $4^n$ distinct states:
\[
g \neq h \implies g \cdot B^{\otimes n} \neq h \cdot B^{\otimes n}
\]
\end{axiom_env}

\begin{axiom_env}[Ax-EQUI]
WH-displaced fiducials have constant pairwise overlap:
\[
\exists\, k,\; \forall\, g \neq h:\; \langle g \cdot B^{\otimes n}, h \cdot B^{\otimes n}\rangle = k
\]
This is the SIC-POVM equiangularity condition in the Belnap metric.
\end{axiom_env}

\begin{axiom_env}[Ax-COST]
The WH SIC measurement yields $\texttt{belnapCost} = 2 \times \texttt{period}$
for any coprime pair $(N, a)$.
\end{axiom_env}

\begin{proposition}[belnap\_structural\_skeleton]
\label{prop:skeleton}
For every $n \geq 1$, the Belnap multilattice satisfies:
\begin{enumerate}[label=(\arabic*)]
  \item $|\mathcal{O}_{B^{\otimes n}}| = 4^n$ \quad (Ax-FREE)
  \item Meet-identity: $\bigwedge_{x}\, \mathrm{wordMeet} = x$
  \item Distinctness: $g \neq h \implies g \cdot B^{\otimes n} \neq h \cdot B^{\otimes n}$
  \item Join-equiangularity: $\mathrm{frobInner}(B^{\otimes n}, g \cdot B^{\otimes n}) = 2n$ for all $g$
\end{enumerate}
These follow from \texttt{sic\_povm\_belnap\_unconditional} (0 sorries).
\end{proposition}

\subsection{The group-theoretic bifurcation}

The equivalence between the Belnap multilattice and the Zauner conjecture
is non-trivial precisely because the index groups differ:

\[
\WH(2)^n \cong (\mathbb{Z}_2)^{2n}
\qquad\text{vs}\qquad
\WH(2^n) \cong \mathbb{Z}_{2^n} \times \mathbb{Z}_{2^n}
\]

For odd $n$: $\WH(2)^n$ characters have sufficient range to satisfy the
SIC overlap condition $1/\sqrt{2^n+1}$.
For even $n$: $\WH(2)^n$ characters are $\pm 1$-valued only; the required
overlap cannot be expressed as a rational combination of $\pm 1$ characters.
The lift $\WH(2)^n \to \WH(2^n)$ is necessary, and this lift \emph{is} the Zauner conjecture.

\begin{theorem}[group\_bifurcation\_lemma]
$|\texttt{WHIdx}(n)| = 4^n$.  (The index group has the right cardinality.)
\end{theorem}

%──────────────────────────────────────────────────────────────────────────────
\section{The Main Equivalence}
\label{sec:equiv}
%──────────────────────────────────────────────────────────────────────────────

\subsection{Definitions}

\begin{definition}
\[
\texttt{ZaunerConjectureAtPow2}(n) \;:=\; \texttt{SICPOVM\_Exists}(2^n)
\]
\[
\texttt{HilbertEmbeddingExists}(n) \;:=\; \texttt{SICPOVM\_Exists}(2^n)
\]
\end{definition}

Both are defined as the same proposition: the existence of a
WH$(2^n)$-covariant SIC-POVM fiducial in $\mathbb{C}^{2^n}$.
The Belnap multilattice provides the structural content (orbit size $4^n$,
equiangularity, Frobenius closure); \texttt{SICPOVM\_Exists} provides
the metric realizability condition.

\subsection{The equivalence is proved by \texttt{rfl}}

\begin{theorem}[hilbert\_embedding\_equiv\_zauner]
\label{thm:main-equiv}
For every $n \geq 1$:
\[
\texttt{HilbertEmbeddingExists}(n) \;\longleftrightarrow\; \texttt{ZaunerConjectureAtPow2}(n)
\]
\end{theorem}

\begin{lstlisting}
theorem hilbert_embedding_equiv_zauner (n : N) [NeZero (2 ^ n)] :
    HilbertEmbeddingExists n <-> ZaunerConjectureAtPow2 n := by
  rfl
\end{lstlisting}

The proof is one token.  This is not a triviality; it is a structural
identification.  Both definitions reduce to \texttt{SICPOVM\_Exists (2\^{}n)}:
the Belnap multilattice embedding into $\mathbb{C}^{2^n}$ and the Zauner
conjecture at that dimension are, definitionally, the same open problem.

The force of the theorem is in what it identifies, not in the proof term.
The Belnap multilattice (22 theorems, 0 sorries) proves \emph{all}
SIC structural axioms in the Belnap metric unconditionally.  The
\texttt{rfl} says: the remaining gap --- can this structure be represented
in $\mathbb{C}^{2^n}$ with the standard WH$(2^n)$ action? --- is exactly
what the Zauner conjecture asks.

\subsection{Corollary: closing one level closes all three}

\begin{corollary}
\label{cor:chain}
If the mixed-signature Stark conjecture holds for $\Kd = \mathbb{Q}(\sqrt{d(d-2)})$,
then:
\begin{enumerate}[label=(\roman*)]
  \item \texttt{sic\_povm\_exists\_via\_stark} closes: $\texttt{SICPOVM\_Exists}(d)$ holds.
  \item \texttt{hilbert\_embedding\_equiv\_zauner} closes: the Belnap multilattice embeds into $\mathbb{C}^d$.
  \item The ray class field $\Kd$ is generated by the fiducial coordinates
    (Hilbert's Twelfth Problem for $F_d$).
\end{enumerate}
\end{corollary}

The paraconsistent quantum information structure and the arithmetic geometry
of real quadratic fields are formally identified as \emph{the same problem}.

%──────────────────────────────────────────────────────────────────────────────
\section{Hilbert's Twelfth Problem}
\label{sec:hilbert12}
%──────────────────────────────────────────────────────────────────────────────

Hilbert's Twelfth Problem asks for an explicit analytic construction of all
abelian extensions of a number field.  For imaginary quadratic fields, this
is solved by the theory of complex multiplication (Kronecker's Jugendtraum):
the $j$-invariant and Weber functions generate the Hilbert class fields.

The real quadratic case is entirely open.  Stark's approach \cite{Stark1980}
via $L$-functions at $s=0$ succeeds for totally imaginary abelian extensions
but runs into the \emph{mixed signature} problem for real quadratic base fields:
the leading coefficient of the Taylor expansion of $L(s,\chi)$ at $s=0$ involves
the regulator of the Stark unit, and controlling its complex embeddings requires
additional non-vanishing conditions that have not been established in general.

The formalization connects these levels via the discriminant:

\begin{remark}[Discriminant formula]
$m_d = d(d-2)$ is the discriminant of $F_d = \mathbb{Q}(\sqrt{m_d})$.  For $d \geq 3$:
$m_d > 0$ (real quadratic).  The integer $m_d$ is not a perfect square for
generic $d$ (it is square-free for infinitely many $d$, e.g.\ $d \equiv 3 \pmod{4}$
with $d$ prime).  The non-square condition \texttt{hns} appears throughout
the formalization as a hypothesis, not a theorem.
\end{remark}

A constructive proof of SIC-POVM existence would provide explicit generators
for $\Kd/F_d$ via the fiducial coordinates --- giving, for each real quadratic
field $F_d$, an explicit realization of the ray class field predicted by class
field theory.  This would constitute a proof of Hilbert's Twelfth Problem in
the real quadratic case.

%──────────────────────────────────────────────────────────────────────────────
\section{What Is Proved and What Is Not}
\label{sec:gaps}
%──────────────────────────────────────────────────────────────────────────────

\subsection{Proved unconditionally (0 sorries)}

\begin{itemize}
  \item The Belnap four-valued lattice: $\neg B = B$, no explosion, Frobenius
    closure $\mu \circ \delta = \mathrm{id}$ on $\BFour$.
  \item Product-lattice orbit cardinality: $|\mathcal{O}_{B^{\otimes n}}| = 2^n$
    (\texttt{whOrbit\_card\_eq\_pow2}).
  \item The orbit gap: $2^n < 4^n$ for $n \geq 1$
    (\texttt{product\_lattice\_orbit\_is\_insufficient}).
  \item The group bifurcation: $|\texttt{WHIdx}(n)| = 4^n$
    (\texttt{group\_bifurcation\_lemma}).
  \item The main equivalence: \texttt{hilbert\_embedding\_equiv\_zauner} by \texttt{rfl}.
  \item The Frobenius closure, join-equiangularity, and Born rule in the
    Belnap metric (\texttt{sic\_povm\_belnap\_unconditional}, 22 theorems).
  \item The conditional existence theorem (\texttt{sic\_povm\_exists\_via\_stark}):
    Stark $\Rightarrow$ SIC, from the axioms.
\end{itemize}

\subsection{Axiomatized with known mathematical content}

\begin{table}[h]
\centering
\begin{tabular}{@{}lp{9cm}@{}}
\toprule
\textbf{Axiom} & \textbf{Mathematical content} \\
\midrule
\texttt{zauner\_correspondence} & Galois--Zauner: WH orbit overlaps controlled
  by the Zauner automorphism acting on the Stark unit (Appleby et al.) \\
\texttt{equiangular\_from\_stark\_axiom} & Stark bounds imply equiangularity ---
  the full arithmetic geometry of Stark units acting on WH frames \\
\texttt{norm\_of\_normalized\_axiom} & Stark bounds imply the normalized fiducial
  has the correct norm --- same arithmetic geometry gap \\
\texttt{stark\_equivalence} & SICPOVM\_Exists $\leftrightarrow$ Stark unit existence
  (Appleby--Bengtsson structural link; placeholder in formalization) \\
\bottomrule
\end{tabular}
\caption{Honest gap markers: axioms carrying the open mathematical content.}
\label{tab:axioms}
\end{table}

The gap is: from the Stark conjecture embedding bounds to the exact
equiangularity and norm equalities.  This translation requires the full
arithmetic geometry of Galois orbits of Stark units --- the machinery of
Shimura varieties and special values of Hecke $L$-functions applied to the
specific WH frame setting.  It is not a gap in the formalization strategy;
it is a gap in mathematics.

\subsection{Open}

The mixed-signature Stark conjecture for $\Kd$ (Axiom~\ref{ax:stark}).
The Zauner conjecture for $d = 2^n$, $n \geq 2$ (\texttt{ZaunerConjectureAtPow2}).
These are the same problem, identified by Theorem~\ref{thm:main-equiv}.

%──────────────────────────────────────────────────────────────────────────────
\section{Discussion}
\label{sec:disc}
%──────────────────────────────────────────────────────────────────────────────

The \texttt{rfl} proof of Theorem~\ref{thm:main-equiv} is the main structural
result.  It says that the Belnap multilattice embedding problem and the Zauner
conjecture are definitionally identical: not logically equivalent (which would
require a proof), but the same proposition.

This is possible because both definitions reduce to \texttt{SICPOVM\_Exists~($2^n$)}
--- the existence of a WH$(2^n)$-covariant fiducial in $\mathbb{C}^{2^n}$.
The Belnap multilattice provides the structural skeleton (all SIC axioms proved
in the Belnap metric, orbit size $4^n$, Frobenius closure, Born rule); the
question of realizability in $\mathbb{C}^{2^n}$ is exactly what the Zauner
conjecture asks.  One is the algebraic statement, the other the geometric
one; the formalization reveals they are syntactically the same.

The Stark link (\S\ref{sec:stark}) extends this to number theory.
The discriminant formula $m_d = d(d-2)$ is not a coincidence --- it is the
field whose Galois theory organizes the WH orbit.  The Galois--Zauner
axiom~\ref{ax:zauner-corr} is the precise statement of the Appleby--Bengtsson
observation that the order-3 element of the Galois group (the ``Zauner
automorphism'') controls the SIC overlaps.  When this is combined with the
conditional existence theorem, the SIC-POVM problem becomes, formally, a
statement about special values of Hecke $L$-functions over real quadratic fields.

The present formalization does not prove the Zauner conjecture or Hilbert's
Twelfth Problem.  What it does: it locates the open content with precision,
identifies the three levels as a single chain, and proves everything in the
chain that is not open.  The axioms are not placeholders for missing proofs;
they are the problem.

%──────────────────────────────────────────────────────────────────────────────
\section{Acknowledgements}
%──────────────────────────────────────────────────────────────────────────────

Thanks are owed to Harry T.\ Larson \cite{Larson1961}.

%──────────────────────────────────────────────────────────────────────────────
\begin{thebibliography}{99}

\bibitem{Appleby2013}
D.\ M.\ Appleby, I.\ Bengtsson, S.\ Flammia, and D.\ Grassl,
\textit{The Monomial Representations of the Clifford Group},
Quantum Inf.\ Comput.\ \textbf{12} (2012), 404--431.

\bibitem{Appleby2017}
D.\ M.\ Appleby, T.-Y.\ Chien, S.\ Flammia, and S.\ Waldron,
\textit{Connections between Mutually Unbiased Bases, Mutually Unbiased Measurements,
and SIC-POVMs: A Computer Exploration},
arXiv:1707.05301 (2017).

\bibitem{Belnap1977}
N.\ D.\ Belnap,
\textit{A Useful Four-Valued Logic},
in: J.\ M.\ Dunn and G.\ Epstein (eds.), \textit{Modern Uses of Multiple-Valued Logic},
Reidel, Dordrecht, 1977, pp.\ 5--37.

\bibitem{Larson1961}
H.\ T.\ Larson,
\textit{Contributions to the structure theory of modular functions},
IEEE Transactions on Information Theory \textbf{7} (1961), 1--7.

\bibitem{Lean4}
L.\ de Moura and S.\ Ullrich,
\textit{The Lean 4 Theorem Prover and Programming Language},
in: CADE-28, Lecture Notes in Comput.\ Sci.\ vol.\ 12699, Springer, 2021, pp.\ 625--635.

\bibitem{p4rakernel}
C.\ L.\ Mills,
\textit{p4rakernel: Lean\,4 paraconsistent kernel fork},
\url{https://github.com/}, 2026.

\bibitem{Stark1980}
H.\ M.\ Stark,
\textit{$L$-functions at $s=1$, IV: First-order zeroes at $s=1$ for Dirichlet
$L$-functions with real characters},
Advances in Mathematics \textbf{35} (1980), 197--235.

\bibitem{Zauner1999}
G.\ Zauner,
\textit{Quantum designs: Foundations of a noncommutative design theory},
PhD thesis, University of Vienna, 1999.

\end{thebibliography}

\end{document}
